\bookchapter{Capitolo 69: Manuale di consultazione rapida e risorse}

\chapterintro{Questo capitolo è un manuale di consultazione rapida utilizzabile in autonomia. Quando non ti ricordi una sintassi o devi decidere in fretta, apri direttamente qui. Conservalo come una “carta da taschino”: codice d'uso comune, domande frequenti, risorse per studiare, senza saltare una pagina.}

\booksection{Scheda rapida del codice d'uso comune}

\begin{center}
\begin{tabularx}{\linewidth}{|X|X|}
\hline
\textbf{Cosa vuoi fare} & \textbf{Codice chiave} \tabularnewline
\hline
Stato privato della pagina & @State private var x = valoreIniziale \tabularnewline
Logica tra più pagine & \texttt{@Observable class VM \{ var x = 0 \}} \tabularnewline
Layout verticale/orizzontale & \texttt{VStack \{ \}}/\texttt{HStack \{ \}} \tabularnewline
Spazio elastico & \texttt{Spacer()} \tabularnewline
Elenco & List \{ ForEach(array) \{ ... \} \} \tabularnewline
Salto di pagina & \texttt{NavigationStack + NavigationLink} \tabularnewline
Finestra dal basso & .sheet(isPresented: \$show) \{ contenuto \} \tabularnewline
Chiudere la pagina corrente & \texttt{@Environment(\textbackslash{}.\allowbreak dismiss) var dismiss} \tabularnewline
Scaricare dati dalla rete & \texttt{try await URLSession.\allowbreak shared.\allowbreak data(from: url)} \tabularnewline
Analisi del JSON & \texttt{struct X: Codable}+\texttt{JSONDecoder().\allowbreak decode} \tabularnewline
Salvare piccole configurazioni & UserDefaults.standard.set(valore, forKey: "chiave") \tabularnewline
Salvare dati strutturati & \texttt{@Model}+\texttt{modelContainer}+\texttt{@Query} \tabularnewline
Adattamento alla modalità scura & Usa i colori di sistema (.primary / Color(.systemGray6)), non valori fissi \tabularnewline
Esegui i test & CmdU \tabularnewline
\hline
\end{tabularx}
\end{center}

\begin{center}\small\color{gray}Scheda di consultazione rapida dello sviluppo iOS (edizione agosto 2026)\end{center}

\booksection{Domande frequenti (FAQ)}

D1: Non ho un Mac, posso imparare lo sviluppo iOS?

Puoi imparare il linguaggio, ma non fare lo sviluppo completo. Swift Playgrounds su iPad/iPhone ti permette di portare a termine tutto il contenuto del Capitolo 61; però Xcode esiste solo su macOS, quindi alla fine serve per forza un Mac. Con un budget limitato puoi nel frattempo affittare un Mac nel cloud o comprare un Mac mini usato come passaggio, e decidere l'investimento dopo aver imparato il linguaggio.

D2: Quanto costa imparare lo sviluppo iOS?

Fase di studio: 0 euro (Xcode è gratuito + test su dispositivo reale gratuiti con l'ID Apple). Fase di pubblicazione: account sviluppatore a 99 dollari l'anno (circa 90 euro). Sommando l'hardware di un Mac, questo è tutto il costo. Rispetto ai ritorni — sull'App Store anche un piccolo strumento ha utenti disposti a pagare — i conti tornano.

D3: Cosa imparare prima, SwiftUI o UIKit?

SwiftUI. Nel 2026 i nuovi progetti, i nuovi tutorial e le nuove posizioni lavorative sono centrati su SwiftUI, e in tutto il libro usiamo quello. UIKit è il framework della generazione precedente, ancora molto usato nei progetti più datati — quando sarai pratico con SwiftUI potrai aggiungere UIKit secondo necessità (Apple offre tutorial sull'interoperabilità) e impararlo in fretta. Non invertire l'ordine.

D4: Facile trovare lavoro nello sviluppo iOS?

Il mercato nel 2026: le posizioni base che chiedono solo di “saper disegnare interfacce” sono molto contese, ma gli sviluppatori che “conoscono Swift + SwiftUI + hanno un'app pubblicata” restano ricercatissimi. Lo stipendio mediano delle posizioni iOS è da sempre tra i più alti tra gli stack tecnologici. Per te che stai studiando, il miglior curriculum è quell'app fatta da te che si trova sullo store.

D5: L'app è stata rifiutata dalla revisione: che fare?

Niente panico, è normalissimo — alla prima presentazione le probabilità di rifiuto non sono basse. Apple di solito indica il motivo in modo molto preciso (quale pagina, quale linea guida): correggi punto per punto, rispondi con una spiegazione in App Store Connect e ripresenta l'app. I tre motivi di rifiuto più comuni: uso dei permessi spiegato male, bug nell'interfaccia, metadati (screenshot/descrizione) che non corrispondono al prodotto. Controllare da soli questi tre punti in anticipo ti risparmia un giro di attesa.

\booksection{Scorciatoie da tastiera di uso frequente in Xcode}

\begin{center}
\begin{tabularx}{\linewidth}{|X|X|}
\hline
\textbf{Azione} & \textbf{Scorciatoia} \tabularnewline
\hline
Esegui / Testa & CmdR / CmdU \tabularnewline
Apri l'anteprima della canvas & OptionCmd \tabularnewline
Apri rapidamente un file & CmdO \tabularnewline
Ricerca globale & CmdF \tabularnewline
Formatta i rientri & CtrlI \tabularnewline
Commenta / Togli il commento & Cmd/ \tabularnewline
Pulisci la cache di build & CmdK \tabularnewline
Vai alla definizione & Cmd + clic \tabularnewline
Rinomina (rinomina sicura in tutto il progetto) & Clic destro Refactor Rename \tabularnewline
\hline
\end{tabularx}
\end{center}

\begin{center}\small\color{gray}Scorciatoie da tastiera di uso frequente in Xcode\end{center}

\begin{tcbinfo}{Consiglio: di fronte ai misteri della compilazione, prima pulisci la cache}
“Ieri funzionava, oggi no”: sei volte su dieci sono cache di build impazzite. Con CmdK pulisci e ricompila: risolve più di metà dei problemi inspiegabili. Se non basta, riavvia Xcode (e anche il Mac); se ancora non basta, “cancella la cartella DerivedData” — è la sequenza finale che gli sviluppatori iOS si passano di bocca in bocca.

\end{tcbinfo}

\booksection{Risorse per lo studio consigliate}

\booksection{Risorse per lo studio consigliate}

\begin{center}
\begin{tabularx}{\linewidth}{|X|X|X|}
\hline
\textbf{Tipo} & \textbf{Risorsa} & \textbf{Perché consigliarla} \tabularnewline
\hline
Ufficiale & developer.apple.com & Documentazione ufficiale Apple, con dentro i tutorial SwiftUI/Swift: la fonte più autorevole \tabularnewline
Ufficiale & SwiftUI Tutorials di Apple & Corso ufficiale passo a passo, massima accoglienza per i principianti \tabularnewline
Ufficiale & Sito ufficiale di Swift (swift.org) & Documentazione ufficiale del linguaggio e playground online \tabularnewline
Community & Hacking with Swift & La raccolta di tutorial gratuiti di Paul Hudson, qualità riconosciuta al vertice \tabularnewline
Community & Stack Overflow (tag swiftui) & Il primo posto dove chiedere aiuto in caso di errori \tabularnewline
Video & Canale ufficiale Apple Developer (video WWDC) & Spiegazioni autorevoli delle novità di ogni anno \tabularnewline
Design & Apple Human Interface Guidelines (HIG) & La bibbia ufficiale per creare interfacce dal "tocco Apple" \tabularnewline
Di supporto & Parte 7: "Programmazione con l'AI in pratica" (Capitoli 37-47) & Il metodo completo per sviluppare per iOS con l'aiuto dell'AI \tabularnewline
\hline
\end{tabularx}
\end{center}

\begin{center}\small\color{gray}Elenco delle risorse per imparare lo sviluppo iOS (agosto 2026)\end{center}

\booksection{Conclusione: dal primo App all'ennesimo App}

Guardando indietro, sei passato da “non sapere cosa sia SwiftUI” a installare Xcode, imparare Swift, disegnare le interfacce con SwiftUI, collegare l'app a database e rete, fino a test, pacchetto e pubblicazione — su questa strada hai già completato il tratto più difficile del percorso da zero a uno.

Il mondo della tecnologia cambia in fretta: i framework si aggiornano, gli strumenti si rinnovano, l'AI si prende in carico sempre più lavoro da “tastiera”. Ma una cosa non passa mai di moda: la capacità di scomporre un'idea vaga in passi chiari e realizzarla passo dopo passo. Gli strumenti sono un amplificatore, il modo di ragionare è la fondamenta. L'amplificatore si può sostituire, la fondamenta no.

Non sei più solo "uno che sa usare il telefono": puoi trasformare le idee in testa in app vere che girano su centinaia di milioni di iPhone. Da oggi crea con la tastiera — questo è il lato romantico dell'essere sviluppatori iOS.

\booksection{Bibliografia}

\begin{enumerate}
  \item Apple Inc. (2026). SwiftUI Framework Documentation. developer.apple.com.
  \item Apple Inc. (2026). Swift Programming Language Guide. swift.org.
  \item Apple Inc. (2026). SwiftData Framework Documentation. developer.apple.com.
  \item Apple Inc. (2026). App Store Review Guidelines. developer.apple.com/app-store/review.
  \item Apple Inc. (2026). Human Interface Guidelines (HIG). developer.apple.com/design.
  \item Apple Inc. (2026). Xcode 27 Release Notes and App Distribution Guide.
\end{enumerate}
